\scanpage{4}
\mspage{2}
\deleted{s'envoie injectivement dans un groupe algébrique
commutatif de rang unipotent nul
(i.e. ext. d'une VA par un tore) alors
\uline{l'application}}
\diagram{$\deleted{X(\overline{\mathbf Q})\longrightarrow J}$}

\textcircled{3} Supposons $X$ de t.f. sur $\overline{\mathbf Q}$. Alors on a
un ``germe de relèvement'' de $\Gamma$
\diagram{$(9)\qquad$
\begin{tikzcd}[ampersand replacement=\&]
\mathcal G_{\overline X}\arrow[r]\&\Gamma\arrow[d,phantom,"\cup" description]\\
\&\Gamma'\arrow[ul]
\end{tikzcd}}

sur ss-groupe ouvert (modulo la relation
de coïncidence sur un ss groupe
ouvert plus petit). Pour \inserted{$\xi\in X(\overline{\mathbf Q})$ donné, et pour}
$\Gamma'$ assez petit,
\inserted{de $\xi$} $\Gamma'
\mathrel{\vcenter{\hbox{\begin{tikzpicture}[x=1mm,y=1mm]
\draw[line width=.45pt] (0,0)--(6,0)--(6,-.4);
\end{tikzpicture}}}}
\mathcal G_{\overline X}$ se factorise par $\mathcal G_{\overline X}(\xi)$, de
\margin{$X=\overline X$}
sorte qu'on a un \inserted{germe d'} hom composé
\[
\begin{gathered}
\hspace{1.6cm}\inserted{\Gamma'\longrightarrow\mathcal E_X}\\[-.2em]
\deleted{\Gamma'\longrightarrow\mathcal E_X}\hspace{1.6cm}
\end{gathered}
\]
rendant commutatif
\diagram{$(10)\qquad$
\begin{tikzcd}[ampersand replacement=\&]
\mathcal E_X\arrow[r]\&\mathcal G_X\arrow[r]\&\Gamma\\
\&\&\Gamma'\arrow[u,hook]\arrow[ull,"{\varphi_\xi^t}"{below}]
\end{tikzcd}}

-- du moins si $\widetilde X=\widetilde X(\xi)$ -- sinon un
hom defini mod composition par un \unsure{autom}
int. provenant de $\pi_1(X)$. \unsure{Soulignons}
\mspage{3}
